\lecture{28}{Консенсус и средний консенсус}

\sourcevideo
    {https://www.youtube.com/playlist?list=PLOzRYVm0a65fhKDS8cYIC7YCuVGJ4FOJx}
    {Week 8: Lecture 28: Consensus and Average Consensus}

Пусть агенты обмениваются состояниями только с соседями по
коммуникационному графу. Консенсус означает сходимость всех состояний
к одному значению, а средний консенсус дополнительно требует, чтобы
это значение совпадало со средним начальных состояний. Ниже
рассматриваются распространение информации по графу, максимальный
консенсус и линейная модель Де Гроота.

\section{Распространение информации по графу}

Рассмотрим ориентированный граф
\[
    \mathcal G=(\mathcal V,\mathcal E),
    \qquad
    \mathcal V=\{1,\ldots,N\}.
\]
Зафиксируем соглашение: \(a_{ij}>0\), если агент \(i\) получает
информацию от агента \(j\). Тогда строка \(i\) матрицы смежности
\(A=(a_{ij})\) описывает входящие связи вершины \(i\), а
\[
    (Ax)_i=\sum_{j=1}^{N}a_{ij}x_j
\]
является взвешенной суммой состояний, доступных за один шаг. При
противоположном соглашении о направлении рёбер входящее
агрегирование записывается через \(A^{\mathsf T}x\).

Элемент степени матрицы имеет вид
\[
    (A^k)_{ij}
    =
    \sum_{i_1,\ldots,i_{k-1}}
    a_{ii_{k-1}}a_{i_{k-1}i_{k-2}}\cdots a_{i_1j}.
\]
Он равен суммарному весу ориентированных блужданий длины \(k\) от
\(j\) к \(i\); в невзвешенном графе это число таких блужданий.
Поэтому при линейном обновлении
\[
    x^{k+1}=Ax^k,
    \qquad
    x^k=A^kx^0,
\]
через \(k\) шагов на состояние вершины могут влиять начальные
значения из её \(k\)-шаговой входящей окрестности. Тот же принцип
лежит в основе линейных графовых слоёв
\[
    H^{\ell+1}
    =
    \sigma\!\left(\widetilde A H^\ell W^\ell\right),
\]
где между операциями агрегирования добавляются обучаемые матрицы,
нормировка и нелинейность.

\section{Консенсус и максимальный консенсус}

\begin{definition}[Консенсус]
Последовательность \(x^k\in\mathbb R^N\) достигает консенсуса, если
существует \(c\in\mathbb R\), для которого
\[
    x^k\longrightarrow c\mathbf1.
\]
\end{definition}

Обычный консенсус не фиксирует значение \(c\). Например, итерация
\(x^{k+1}=\tfrac12x^k\) приводит все состояния к нулю, но не сохраняет
информацию о начальных данных.

\begin{definition}[Средний консенсус]
Алгоритм достигает среднего консенсуса, если
\[
    x^k
    \longrightarrow
    \frac1N\mathbf1\mathbf1^{\mathsf T}x^0.
\]
\end{definition}

Для связного неориентированного графа максимальный консенсус задаётся
обновлением
\[
    x_i^{k+1}
    =
    \max\bigl\{x_i^k,\,x_j^k:j\in\mathcal N_i\bigr\}.
\]
Каждая последовательность \(x_i^k\) не убывает и не превосходит
начальный глобальный максимум. Если максимум первоначально находится
в вершине \(j^\star\), то после \(k\) шагов его знают все вершины на
расстоянии не более \(k\) от \(j^\star\).

\begin{theorem}[Конечное время максимального консенсуса]
На связном неориентированном графе
\[
    x_i^k=\max_j x_j^0
\]
для всех \(i\) и всех
\[
    k\ge\operatorname{diam}(\mathcal G).
\]
\end{theorem}

Аналогичное обновление с минимумом достигает минимального консенсуса.
Среднее устроено сложнее: необходимо одновременно распространить
информацию, сохранить её суммарную массу и устранить рассогласование.
Передача всех начальных значений позволяет получить точный ответ за
конечное время, но требует памяти и сообщений, растущих с \(N\).
Линейные алгоритмы используют сообщения фиксированной размерности и
обычно сходятся лишь асимптотически.

\section{Модель Де Гроота и стохастические матрицы}

В модели Де Гроота агент обновляет состояние по правилу
\[
    x_i^{k+1}=\sum_{j=1}^{N}w_{ij}x_j^k,
    \qquad
    x^{k+1}=Wx^k.
\]
Матрица \(W\) согласована с графом: \(w_{ij}=0\), если агент \(i\) не
получает сообщение от \(j\).

\begin{definition}[Строчная стохастичность]
Неотрицательная матрица \(W\) называется строчно-стохастической, если
\[
    W\mathbf1=\mathbf1.
\]
\end{definition}

Тогда каждое новое состояние является выпуклой комбинацией
предыдущих, поэтому
\[
    \min_jx_j^k
    \le x_i^{k+1}\le
    \max_jx_j^k.
\]
Однако строчной стохастичности недостаточно для сходимости. Матрица
\[
    W=
    \begin{pmatrix}
        0&1\\
        1&0
    \end{pmatrix}
\]
дважды стохастична, но \(W^2=I\), поэтому состояния периодически
меняются местами.

\begin{definition}[Примитивность]
Неотрицательная матрица \(W\) называется примитивной, если существует
\(m\ge1\), для которого \(W^m>0\) покомпонентно.
\end{definition}

Для стохастической матрицы примитивность исключает разложение сети и
периодические колебания. В графовых терминах она соответствует
сильной связности и апериодичности; на связном неориентированном графе
апериодичность обычно обеспечивается положительными диагональными
весами.

\begin{theorem}[Предел степеней]
Пусть \(W\) --- примитивная строчно-стохастическая матрица. Тогда
существует единственный вектор \(\pi>0\), удовлетворяющий
\[
    \pi^{\mathsf T}W=\pi^{\mathsf T},
    \qquad
    \mathbf1^{\mathsf T}\pi=1,
\]
и
\[
    W^k\longrightarrow\mathbf1\pi^{\mathsf T}.
\]
Следовательно,
\[
    x^k\longrightarrow
    \mathbf1\pi^{\mathsf T}x^0.
\]
\end{theorem}

Предел является взвешенным средним начальных состояний. Вес агента
\(i\) равен \(\pi_i\).

\section{Средний консенсус}

\begin{definition}[Двойная стохастичность]
Неотрицательная матрица \(W\) называется дважды стохастической, если
\[
    W\mathbf1=\mathbf1,
    \qquad
    \mathbf1^{\mathsf T}W=\mathbf1^{\mathsf T}.
\]
\end{definition}

Столбцовая стохастичность сохраняет сумму состояний:
\[
    \mathbf1^{\mathsf T}x^{k+1}
    =
    \mathbf1^{\mathsf T}Wx^k
    =
    \mathbf1^{\mathsf T}x^k.
\]
Если одновременно достигается консенсус \(x^k\to c\mathbf1\), то
\[
    Nc=\mathbf1^{\mathsf T}x^0,
\]
то есть \(c\) равно начальному среднему.

\begin{theorem}[Средний консенсус]
Если \(W\) примитивна и дважды стохастична, то
\[
    W^k
    \longrightarrow
    J
    :=
    \frac1N\mathbf1\mathbf1^{\mathsf T},
\]
и для любого \(x^0\)
\[
    x^k\longrightarrow Jx^0.
\]
\end{theorem}

\begin{proof}
Примитивность и строчная стохастичность дают
\(W^k\to\mathbf1\pi^{\mathsf T}\). По столбцовой стохастичности
\(N^{-1}\mathbf1\) является нормированным левым собственным вектором
для собственного значения \(1\). Единственность стационарного вектора
даёт \(\pi=N^{-1}\mathbf1\).
\end{proof}

Двойная стохастичность без примитивности недостаточна, как показывает
матрица перестановки. Если \(W=W^{\mathsf T}\) и
\(W\mathbf1=\mathbf1\), то после транспонирования получаем
\(\mathbf1^{\mathsf T}W=\mathbf1^{\mathsf T}\); поэтому симметричная
строчно-стохастическая матрица автоматически дважды стохастична.

Проекторы
\[
    J=\frac1N\mathbf1\mathbf1^{\mathsf T},
    \qquad
    J_\perp=I-J
\]
выделяют консенсусную часть \(Jx=\bar x\mathbf1\) и рассогласование
\(J_\perp x=x-\bar x\mathbf1\). Для дважды стохастической матрицы
\(WJ=JW=J\), поэтому
\[
    J_\perp x^{k+1}
    =
    (W-J)J_\perp x^k.
\]
Средний консенсус достигается тогда, когда оператор \(W-J\) сжимает
подпространство рассогласования.

\section{Спектральная скорость и лапласовский шаг}

Если \(W\) диагонализируема,
\[
    W=P\Lambda P^{-1},
    \qquad
    W^k=P\Lambda^kP^{-1}.
\]
Поэтому при простом собственном значении \(1\) и условии
\[
    |\lambda_i(W)|<1,
    \qquad i\ge2,
\]
степени матрицы сходятся к проектору на консенсусное подпространство.
Для вещественной симметричной дважды стохастической матрицы введём
\[
    \rho_\perp(W)
    =
    \max_{i\ge2}|\lambda_i(W)|.
\]
Тогда
\[
    \lVert W^k-J\rVert_2
    =
    \rho_\perp(W)^k
\]
и
\[
    \lVert x^k-Jx^0\rVert
    \le
    \rho_\perp(W)^k
    \lVert x^0-Jx^0\rVert.
\]
Чем меньше \(\rho_\perp(W)\), тем быстрее затухает рассогласование.

Для связного неориентированного графа можно взять
\[
    W=I-\alpha L,
\]
где \(L\) --- матрица Лапласа. Локальное обновление имеет вид
\[
    x_i^{k+1}
    =
    x_i^k
    -
    \alpha
    \sum_{j\in\mathcal N_i}
    a_{ij}(x_i^k-x_j^k).
\]
Поскольку \(L\mathbf1=0\) и \(L=L^{\mathsf T}\), матрица \(W\)
сохраняет среднее. Для неотрицательности весов достаточно, например,
\[
    0<\alpha\le\frac1{d_{\max}}
\]
в невзвешенном графе. Если
\[
    0=\lambda_1(L)<\lambda_2(L)\le\cdots\le\lambda_N(L),
\]
то собственные значения \(W\) равны
\[
    1-\alpha\lambda_i(L).
\]
Спектральная сходимость обеспечивается условием
\[
    0<\alpha<\frac{2}{\lambda_N(L)},
\]
а её скорость определяется величиной
\[
    \max_{i\ge2}|1-\alpha\lambda_i(L)|.
\]
В частности, малое значение алгебраической связности
\(\lambda_2(L)\) замедляет консенсус.

Если матрица не диагонализируема, используется жорданова форма.
Степень блока \(J_\lambda=\lambda I+N\) размера \(r\) равна
\[
    J_\lambda^k
    =
    \sum_{s=0}^{r-1}
    \binom{k}{s}\lambda^{k-s}N^s.
\]
При \(|\lambda|<1\) экспоненциальное убывание подавляет
полиномиальные множители. Для сходимости собственное значение \(1\)
должно быть полупростым; у примитивной стохастической матрицы оно
простое, а остальные собственные значения лежат внутри единичного
круга.

\section{Ориентированные и векторные модели}

На ориентированном графе локально естественно строить
строчно-стохастическую матрицу. При сильной связности и
апериодичности она даёт
\[
    x^k\longrightarrow
    \mathbf1\pi^{\mathsf T}x^0,
\]
но в общем случае \(\pi\ne N^{-1}\mathbf1\), поэтому предел является
взвешенным средним. Для устранения этого смещения применяют, например,
push-sum и алгоритмы отношения масс.

Если состояние агента векторное,
\[
    x_i^k\in\mathbb R^d,
\]
то объединённое обновление записывается как
\[
    \mathbf x^{k+1}
    =
    (W\otimes I_d)\mathbf x^k.
\]
При \(W^k\to J\) каждая координата усредняется независимо:
\[
    \mathbf x^k
    \longrightarrow
    \mathbf1\otimes
    \left(
        \frac1N\sum_{i=1}^{N}x_i^0
    \right).
\]

В одном синхронном раунде агент отправляет соседям текущее состояние,
получает их состояния и вычисляет взвешенную сумму. Для векторов
размерности \(d\) суммарный объём передачи по активным рёбрам имеет
порядок
\[
    O(|\mathcal E|d),
\]
а число раундов до заданной точности определяется спектральным
коэффициентом смешивания. Базовый анализ предполагает фиксированный
граф, синхронные итерации, точную передачу сообщений и отсутствие
злонамеренных агентов; задержки, потери пакетов, асинхронность и
изменяющаяся топология требуют отдельных условий.
